% Memo, round 2: three storylines for the Introduction (5 November 2026).
% Draft paragraphs are \input from paragraphs/, so the memo and the three
% Introduction drafts always show the same text. Round 1 is in git history.
% Compile: pdflatex -> bibtex -> pdflatex x2 (from this directory).
\documentclass[11pt]{article}
\usepackage[T1]{fontenc}
\usepackage[a4paper,margin=2.2cm]{geometry}
\usepackage{times}
\usepackage[superscript,biblabel,nomove]{cite}
\usepackage{booktabs,tabularx,array,enumitem}
\usepackage[dvipsnames]{xcolor}
\usepackage[most]{tcolorbox}
\usepackage[hidelinks]{hyperref}
\newcommand{\JournalTitle}[1]{#1}
\setlength{\parskip}{0.45em}
\setlength{\parindent}{0pt}
\setlist{itemsep=0.15em,topsep=0.25em}
\definecolor{ink}{HTML}{23373B}
\definecolor{accent}{HTML}{A65A10}
\newtcolorbox{draft}[1]{colback=ink!3,colframe=ink!45,boxrule=0.4pt,arc=1.5pt,
  left=6pt,right=6pt,top=4pt,bottom=4pt,fonttitle=\small\bfseries,coltitle=ink,
  colbacktitle=ink!10,title={#1}}
\newtcolorbox{keypoint}{colback=accent!6,colframe=accent!70,boxrule=0.5pt,arc=1.5pt,
  left=6pt,right=6pt,top=4pt,bottom=4pt}
\newcolumntype{L}{>{\raggedright\arraybackslash}X}
\newcommand{\old}[1]{\textit{``#1''}}
\title{\large\bfseries Three storylines for the Introduction (round 2)\\[0.2em]
\normalsize Transition, revisions to Storylines 1 and 2, and a Storyline 3 for a multidisciplinary journal}
\date{5 November 2026}
\begin{document}
\maketitle
\vspace{-1.5em}

\begin{keypoint}
\textbf{Summary.}
\begin{itemize}[leftmargin=1.2em]
\item \textbf{Transition.} Paragraph~2 now ends by naming two routes out of the punishment dilemma. Observers can be brought to agree, or reputations can be made robust to a single disputed judgement. Each storyline's paragraph~3 opens by picking up one of the two routes.
\item \textbf{Storyline 1.} The two enumerated questions and the sentence that justified them are replaced by one gap sentence. Paragraph~4 now opens with the concept itself (inertia that can differ between the two directions) instead of \old{rather than instantaneously}.
\item \textbf{Storyline 2.} The gap is reframed from \old{one observation at a time} to a sharper point. When evidence arrives piecemeal, agreement depends on how individuals revise their judgements, and it is useful only if it is also accurate. The critical list of conditions is replaced by a neutral description and your summary sentence.
\item \textbf{Storyline 3 (new).} This version shifts the frame \emph{from information to judgement}. Existing remedies improve what observers know, but a problem of judgement remains: how a standing verdict should be revised as evidence arrives, balancing the condemnation of cooperators against the excusing of defectors. It contains Storyline~1 as its core and Storyline~2 as a consequence.
\item \textbf{Recommendation.} Use Storyline~3 for a multidisciplinary journal (\textit{Nature}, \textit{Nature Human Behaviour}). Use Storyline~1 for \textit{PNAS} or \textit{Nature Communications}. In either case, Storyline~2 enters as one sentence in paragraph~5 and one Discussion paragraph (drafts below).
\end{itemize}
\end{keypoint}

% =============================================================================
\section{The transition at the end of paragraph 2}

\textbf{Diagnosis of round 1.} Paragraph~2 ended with \old{A remedy must therefore keep justified punishers in good standing without letting genuine defectors escape condemnation.} That sentence states a goal but does not point to a lever. The next paragraph therefore had to change topic. In Storyline~1 it began \old{Part of this fragility lies in how reputations respond to evidence}, where ``this fragility'' had no antecedent in the sentence before. In Storyline~2 it began with remedies, which the goal sentence had not prepared.

\textbf{New ending.} The dilemma arises because a justified refusal is judged by an observer who disagrees, and because that single judgement overturns the punisher's standing. There are therefore two levers: remove the disagreement, or remove the single judgement's power to overturn a reputation. The new ending names both and adds the constraint that defectors must still be recognized. That constraint motivates Fig.~5 without a separate sentence later.

\begin{draft}{Paragraph 2 (all storylines), 180 words}
% Paragraph 2 (all storylines): the punishment dilemma.
% Round 2: ends with two routes out of the dilemma, which the storylines pick up.
The difficulty is sharpest for the norms that sustain cooperation under public information. Under these norms, a discriminator who refuses to help an individual in bad standing keeps its own good standing, because the refusal counts as justified punishment\cite{panchanathan2003,ohtsuki2004,okada2020dilemmas}. This exemption works only if the observer shares the donor's view of the recipient, which under private assessment is often not the case. An observer who still regards the recipient as good sees an unjustified refusal and condemns the discriminator. Disagreement about one individual thereby becomes disagreement about those who act towards it, and it spreads through the population\cite{hilbe2018,uchida2013,radzvilavicius2019}. Discriminators thus face a punishment dilemma. Withholding help from defectors protects cooperation against exploitation, yet each act of punishment risks their own reputation in the eyes of observers who disagree\cite{okada2020dilemmas,michelmata2024,pires2025}. The dilemma can be eased in two ways. Observers can be brought to agree about each other's standing, so that justified punishment is recognized as such, or reputations can be made robust enough that a single disputed judgement does not overturn them. In both cases, genuine defectors must still lose their standing.

\end{draft}

\begin{center}\small
\begin{tabularx}{\linewidth}{@{}l L@{}}
\toprule
Storyline & How paragraph 3 picks up the two routes \\
\midrule
1 & One sentence acknowledges the first route, then \emph{``The second route depends on how reputations respond to evidence.''} \\
2 & \emph{``Most remedies take the first route and align what observers know.''} The paragraph then argues that agreement can also come from private revision. \\
3 & \emph{``Existing remedies pursue both routes mainly by acting on the information that observers hold.''} The paragraph then argues that a problem of judgement remains. \\
\bottomrule
\end{tabularx}
\end{center}

% =============================================================================
\section{Storyline 1: revisions}

\subsection*{Are the two explicit questions necessary?}
No. Enumerated questions help when the Results are organized around them, they are short, and nothing else repeats them. In round~1 all three conditions failed.
\begin{itemize}[leftmargin=1.2em]
\item The questions were followed by a sentence justifying them.
\item The same content reappeared in paragraph~4, as the two design choices, and in paragraph~5, as the two findings.
\end{itemize}
Each idea therefore appeared three times. The revision ends paragraph~3 with one declarative sentence that names the two features and links them back to paragraph~2: \emph{``How the direction of reputational change and the treatment of incomplete evidence affect the punishment dilemma, and with it the evolution of cooperation, remains unclear.''}

If you want a hook, one question can replace that sentence: \emph{``Should losing a reputation require as much evidence as regaining it?''} Use one question, not two, and drop the justification.

\subsection*{The redundant sentence}
\old{Both bear directly on the punishment dilemma, because buffering negative evidence protects justified punishers but may also shelter those who defect deliberately\dots} is removed. Its two halves now sit where they belong. The link to the dilemma is the new gap sentence. The risk of sheltering defectors is the last sentence of paragraph~2 (\emph{``In both cases, genuine defectors must still lose their standing''}) and the opportunism finding in paragraph~5.

\subsection*{The opening of paragraph 4}
\old{rather than instantaneously} has three problems.
\begin{enumerate}[leftmargin=1.5em]
\item It defines the model by negation, so the reader learns what the model is not.
\item It implies that the contrast is with instantaneous models. Scores, records and aggregation are also not instantaneous, which makes the novelty sound like accumulation, a point paragraph~3 has just conceded.
\item It delays the key concept.
\end{enumerate}
The new opening names the concept and its distinctive feature at once: \emph{``reputations have inertia, and \dots\ this inertia can differ between losing and regaining standing.''}

\subsection*{Storyline 2 as a supplementary perspective}
One sentence in paragraph~5 carries it: \emph{``Because observers share no information, this agreement emerges from how each individual revises its own judgements.''} A Discussion paragraph develops it (end of Section~3).

\begin{draft}{Storyline 1, paragraph 3, 195 words}
% Storyline 1, paragraph 3 (round 2): the second route and how reputations change.
Much theoretical work has pursued the first route, through public institutions, gossip, empathetic evaluation or the synchronization of private opinions\cite{radzvilavicius2021,kawakatsu2024,radzvilavicius2019,murase2024synchronization,tham2026}. The second route depends on how reputations respond to evidence. In the standard framework, each new assessment replaces an individual's image, so a single disputed act of punishment overturns a standing built over many interactions\cite{ohtsuki2004,hilbe2018}. Models that draw on more evidence before assigning a reputation, such as quantitative scores, simple records, graded reputations or the aggregation of several observations, can make cooperation more robust\cite{schmid2023,clark2020,tanabe2013,murase2022,michelmata2024,yamamoto2025}. Most of them, however, treat reputational change as a single symmetric process, in which positive and negative evidence enter through the same score, threshold or observation window and affect treatment only once a threshold is crossed. Human evaluation departs from both features. Impressions of others are revised asymmetrically, although the direction of the asymmetry depends on context\cite{skowronski1989,slovic1993,mendesiedlecki2013diagnostic,siegel2018}. Institutions, too, grade their response and treat the two directions differently, escalating sanctions only with repeated violations\cite{ostrom1990,cox2010} and restoring an offender's standing only after years without a further offence\cite{blumstein2009}. How the direction of reputational change and the treatment of incomplete evidence affect the punishment dilemma, and with it the evolution of cooperation, remains unclear.

\end{draft}
\begin{draft}{Storyline 1, paragraph 4, 147 words}
% Storyline 1, paragraph 4 (round 2): approach, opening with the positive concept.
Here, we develop a model of indirect reciprocity in which reputations have inertia, and in which this inertia can differ between losing and regaining standing (Fig.~1). An established reputation changes only after a run of contrary assessments. Negative assessments move an individual in good standing through a sequence of warning states before the reputation turns bad, and positive assessments move an individual in bad standing through recovery states before it turns good. The two paths can differ in length, so that losing and regaining standing may require different amounts of evidence. Each intermediate state is interpreted as good or bad with a probability that sets how strongly incomplete evidence already affects treatment. The model thus treats two properties of reputational change as separate design choices, namely how much evidence is needed to change a reputation in each direction and how incomplete evidence is expressed in the meantime.

\end{draft}
\begin{draft}{Storyline 1, paragraph 5, 182 words}
% Storyline 1, paragraph 5 (round 2): findings; one sentence carries the Storyline 2 perspective.
We find that the allocation of inertia between the two directions, rather than its total amount, determines whether private observers agree with one another and still distinguish cooperators from defectors. Because observers share no information, this agreement emerges from how each individual revises its own judgements. The beneficial direction depends on the norm. Compact architectures, selected solely by these reputational criteria, make cooperation stable and easy to reach under each of four classical norms (Scoring, Shunning, Stern Judging and Simple Standing). At the same depth, inertia placed on the wrong path can leave no cooperation at all. We further show that the tolerance that protects justified punishers can also shelter opportunists who defect only occasionally, a vulnerability that the aggregation and score-based rules we test share. Allowing warning states to carry partial consequences closes this loophole for the opportunistic strategies we consider, while cooperation remains high. Reputation-based cooperation therefore depends not only on which behaviours a society judges as good, but also on how readily it revises those judgements in each direction and on how it treats individuals before a judgement changes.

\end{draft}

% =============================================================================
\section{Storyline 2: diagnosis and a refined entry point}

\subsection*{Why the round-1 gap did not work}
The round-1 gap read \old{Can agreement arise from how individuals integrate the evidence they already receive, one observation at a time?} It has three weaknesses.
\begin{enumerate}[leftmargin=1.5em]
\item \textbf{It is defined against a technical detail} of Michel-Mata et al.\cite{michelmata2024}, the number of observations per update. ``One observation at a time'' also sounds like a limitation rather than a strength.
\item \textbf{It is not new as stated.} Quantitative scores and simple records already integrate evidence one observation at a time\cite{schmid2023,clark2020}. A referee would point this out.
\item \textbf{It anticipates its own answer.} \old{In principle it can} answers the question it has just posed.
\end{enumerate}

The two closing sentences had flaws on closer inspection, as you suspected.
\begin{itemize}[leftmargin=1.2em]
\item \old{Independent observers agree whenever each reaches a reliable verdict on its own} is true but hides the real difficulty. Under the punishment dilemma, the assessments a cooperator receives depend on what observers think of its recipients. Disagreement is therefore endogenous and self-reinforcing. Integrating more evidence produces reliable verdicts only if cooperators and defectors receive clearly different rates of positive assessment. That separation is exactly what the direction of inertia controls. The argument is sound for the Discussion, where it can be stated with this qualification, but it is too quick for the Introduction.
\item \old{A rule that makes verdicts reliable by discounting contrary evidence may also discount genuine defection} is generic. It also misses half the point, because inertia can delay recovery as well as loss.
\end{itemize}

The list of conditions (institutions need adherence, gossip must be frequent, empathy needs knowledge of the donor's view) read as a critique of other groups' mechanisms. Your summary sentence makes the same point neutrally. The revision describes what each mechanism does and keeps only the summary.

\subsection*{What was missing: why our model is needed}
The round-1 text never said why the existing private-information routes do not settle the question. The revision makes three steps explicit.
\begin{enumerate}[leftmargin=1.5em]
\item \textbf{Piecemeal evidence.} In many settings evidence about a person arrives one encounter at a time, and observers carry a judgement between encounters. Agreement then depends on the \emph{revision rule}, not only on the amount of information. Michel-Mata et al.\ premise that information is abundant. Our premise is that even abundant information reaches each observer over time.
\item \textbf{Agreement must also be accurate.} Emergent agreement can be wrong: observers who regard almost everyone as bad agree perfectly. This conceptual point needs no citation, and our Fig.~2 demonstrates it.
\item \textbf{Direction.} The open question is which properties of revision make agreement accurate, in particular whether losing and regaining standing should follow the same rule. No existing revision rule has this property, and it is what our results are about.
\end{enumerate}

\subsection*{Opening of paragraph 4}
\old{\dots in which each observer integrates its own observations over time, without sharing information and without observing more than in standard private-assessment models} has the same problem as Storyline~1's opening. Its two ``without'' clauses define the model by what it lacks, and they sound defensive. ``Integrates over time'' is also true of scores. The revision opens with the concept, inertia that can differ between the two directions, and states the informational setting in a separate plain sentence.

\subsection*{Assessment}
Even refined, Storyline~2's distinctive content converges on direction, the content of Storyline~1. Its frame, agreement without shared information, is Michel-Mata et al.'s frame. It is therefore strongest as a supplementary perspective: the sentence in paragraph~5 and the Discussion paragraph below. Use it as a full Introduction only if the target audience is primarily interested in decentralized consensus.

\begin{draft}{Storyline 2, paragraph 3, 176 words}
% Storyline 2, paragraph 3: from information synchronization to individual integration of evidence.
Most proposed remedies keep punishers in good standing by aligning what observers know, because cooperation is stable only when individual opinions are sufficiently correlated\cite{murase2024synchronization}. Public institutions broadcast a single verdict\cite{radzvilavicius2021}, gossip spreads opinions between peers\cite{sommerfeld2007,kawakatsu2024}, and other rules synchronize private opinions directly\cite{murase2024synchronization} or evaluate actions from the donor's perspective\cite{radzvilavicius2019}. Each route rests on conditions that are often unmet. Institutions need monitoring capacity and the adherence of those they serve\cite{radzvilavicius2021}, gossip must be frequent and accurate enough to produce consensus\cite{kawakatsu2024,seki2016}, and empathetic evaluation requires knowing how the donor sees the recipient\cite{radzvilavicius2019}. These conditions are strained where empathy, communication and institutional trust are compromised, as in polarized societies\cite{iyengar2019,michelmata2024}. Recent models therefore seek agreement from private information alone, by letting each observer judge from several fresh observations of the same individual\cite{michelmata2024} or assess recipients as well as donors\cite{tham2026}. Both approaches enlarge the evidence behind each judgement. Whether agreement can also arise from how individuals integrate the evidence they already receive, one observation at a time, remains open. In principle it can, because independent observers agree whenever each reaches a reliable verdict on its own. A rule that makes verdicts reliable by discounting contrary evidence, however, may also discount genuine defection.

\end{draft}
\begin{draft}{Storyline 2, paragraph 4, 113 words}
% Storyline 2, paragraph 4 (round 2): approach, opening with the positive concept.
Here, we develop a model of indirect reciprocity in which each observer revises its private judgements with inertia, and in which this inertia can differ between losing and regaining standing (Fig.~1). Observers share no information, and each update draws on a single observation. An established judgement changes only after a run of contrary assessments. Negative assessments move an individual in good standing through a sequence of warning states before the observer's verdict turns bad, and positive assessments move an individual in bad standing through recovery states before it turns good. Each intermediate state is interpreted as good or bad with a probability that sets how strongly incomplete evidence already affects the observer's behaviour.

\end{draft}
\begin{draft}{Storyline 2, paragraph 5, 174 words}
% Storyline 2, paragraph 5: findings and implication.
We find that private observers can agree with one another and still distinguish cooperators from defectors, although they share no information, provided that inertia is allocated in a norm-specific direction. Agreement alone is not sufficient. With the same small inertia in both directions, observers under Shunning agree only because they regard nearly everyone as bad. Compact architectures, selected solely by these two reputational criteria, make cooperation stable and easy to reach under each of four classical norms (Scoring, Shunning, Stern Judging and Simple Standing). At the same depth, inertia placed on the wrong path can leave no cooperation at all, so how evidence is processed matters more than how much of it is retained. Finally, the tolerance that makes private verdicts robust can also shelter opportunists who defect only occasionally, a vulnerability that the aggregation and score-based rules we test share. Allowing warning states to carry partial consequences closes this loophole for the opportunistic strategies we consider, while cooperation remains high. Consensus about reputations therefore need not be imposed through shared information. It can emerge from how each individual weighs its own evidence.

\end{draft}
\begin{draft}{Storyline 2 as a Discussion paragraph inside Storyline 1 or 3, 148 words}
% Discussion paragraph carrying the Storyline 2 perspective inside Storylines 1 or 3.
Our results also bear on how agreement about reputations arises. Agreement is usually attributed to mechanisms that align information across observers, such as institutions, gossip and opinion synchronization\cite{radzvilavicius2021,kawakatsu2024,murase2024synchronization}, or to observers who draw on several observations at once\cite{michelmata2024}. In our model, observers share nothing and update on one observation at a time, yet they largely agree. In the mean-field theory, agreement arises because inertia makes each observer's verdict a reliable function of the target's behaviour, and independent observers who each reach a reliable verdict agree. This route complements information sharing rather than replacing it. It remains available where shared information is costly or untrusted, and it can be combined with institutions or gossip where they exist. Emergent agreement is not automatically useful, however. Under one-step inertia, observers using Shunning agree because they regard nearly everyone as bad. The direction of inertia decides whether emergent agreement is also accurate.

\end{draft}

% =============================================================================
\section{Storyline 3: from information to judgement}

\subsection*{What a multidisciplinary journal needs}
Editors of \textit{Nature} and similar journals look for four things.
\begin{enumerate}[leftmargin=1.5em]
\item A problem recognizable outside the field.
\item A reframing that changes how the field thinks, rather than a better mechanism within the current frame.
\item A simple idea with a sharp result.
\item Implications that travel, for example to institutions, online platforms and artificial agents.
\end{enumerate}
Storyline~1 delivers the third point well but stays within the theory of indirect reciprocity. Storyline~2 has broad appeal, but its frame belongs to Michel-Mata et al. Storyline~3 aims at all four.

\subsection*{Thesis}
The field has treated the failure of private reputations as a problem of \emph{information}: observers know too little, or know different things. Better information reduces disputed judgements but cannot remove them, because observation and assessment stay noisy. What remains is a problem of \emph{judgement}: how a standing verdict about a person should be revised as evidence arrives. Every revision rule trades the condemnation of cooperators on disputed evidence against the excusing of defectors whose misconduct is not yet conclusive. Societies manage this trade-off through procedure, but models revise reputations by a single symmetric rule.

\subsection*{Logic chain}
\begin{enumerate}[leftmargin=1.5em]
\item The punishment dilemma has two routes out (paragraph~2).
\item Existing remedies pursue both routes by acting on information, making it shared (institutions, gossip, synchronization) or richer (scores, records, aggregation).
\item Information cannot remove noise, so a judgement problem remains: how to revise a standing verdict.
\item This is a trade-off between two errors, condemning cooperators and excusing defectors. The trade-off is familiar far beyond indirect reciprocity, from legal procedure to content moderation.
\item Institutions handle it procedurally and asymmetrically: warnings, escalating sanctions, slow restoration of standing. Human impressions are also revised asymmetrically.
\item Models use one symmetric rule. How the revision rule shapes the dilemma, and how it should be designed, is open.
\item Our model has two adjustable features: direction (depths) and partial consequences (strengths). Observers share nothing, so any agreement emerges from the rule. This carries Storyline~2's insight.
\item Results:
  \begin{itemize}[leftmargin=1.2em]
  \item Direction, not amount, decides accuracy and agreement.
  \item The right direction mirrors the norm's typical misjudgement.
  \item The wrong direction can eliminate cooperation.
  \item The most protective rule shelters opportunists, and partial consequences restore deterrence.
  \end{itemize}
\end{enumerate}

\subsection*{Why it fits, and what it demands}
\begin{itemize}[leftmargin=1.2em]
\item \textbf{Reframing.} ``What observers know'' versus ``how they revise what they believe'' is a clean conceptual contrast that the closing sentence states for a general reader.
\item \textbf{Fit to results.}
  \begin{itemize}[leftmargin=1.2em]
  \item The trade-off between two errors gives direction (Figs.~2--4) and partial consequences (Fig.~5) a single motivation.
  \item The ``mirrors the misjudgement'' finding comes from the Discussion and makes the norm dependence intuitive rather than a list of cases.
  \item Fig.~5a has the strongest role here: aggregation and scores, which are information-rich rules, share the opportunism window. This supports the claim that the remaining problem lies in judgement rather than information.
  \end{itemize}
\item \textbf{Demands.} The Discussion must carry the broad claims. It needs a paragraph on procedures in real reputation systems and a clear statement of scope: one family of opportunists, mean-field theory, fixed architectures.
\item \textbf{Risks.} Overreach about human procedures. Keep the institutional examples cited and few, and use ``often'' and ``mostly'' as in the draft. The word ``procedure'' should not suggest a legal analysis. In the paper it means the revision rule.
\item \textbf{Further examples (optional, need sources).} Penalty-point systems for driving licences, in which points accumulate to a suspension and expire after a period, and the warning-and-strike policies of online platforms. Both have the same structure as our model: partial consequences before the verdict, and separate rules for losing and regaining standing. I could not check sources from here. If used, cite official policy documents.
\end{itemize}

\subsection*{Optional broader first sentence for paragraph 1}
For a multidisciplinary journal, the opening can start from everyday experience before the theory: \emph{``From village communities to online marketplaces, people decide whom to help by what they know of others' past conduct\cite{milinski2016,tadelis2016}.''} The rest of paragraph~1 stays as it is.

\begin{draft}{Storyline 3, paragraph 3, 185 words}
% Storyline 3, paragraph 3: from information to judgement.
Existing remedies pursue both routes mainly by acting on the information that observers hold. They make it shared, through public institutions, gossip or the synchronization of private opinions\cite{radzvilavicius2021,kawakatsu2024,murase2024synchronization,tham2026}, or richer, through scores, records or the aggregation of several observations\cite{schmid2023,clark2020,michelmata2024}. Better information reduces erroneous and disputed judgements but does not eliminate them, because observation and assessment remain noisy\cite{hilbe2018,fujimoto2023}. What remains is a problem of judgement, namely how a standing verdict about another person should be revised as new evidence arrives. Any revision rule faces a trade-off between condemning cooperators on disputed evidence and excusing defectors whose misconduct is not yet conclusive. Human institutions often manage this trade-off through procedure. They escalate sanctions only with repeated violations\cite{ostrom1990,cox2010} and restore an offender's standing only after a period of good conduct\cite{blumstein2009}, and impressions of other people are likewise revised asymmetrically\cite{skowronski1989,slovic1993,mendesiedlecki2013diagnostic,siegel2018}. Models of indirect reciprocity, by contrast, mostly revise reputations through a single symmetric rule, whether the latest assessment, a score or a fixed window of observations\cite{ohtsuki2004,hilbe2018,schmid2023,michelmata2024}. How the rule by which private judgements are revised shapes the punishment dilemma, and how such a rule should be designed, remains open.

\end{draft}
\begin{draft}{Storyline 3, paragraph 4, 113 words}
% Storyline 3, paragraph 4: approach.
Here, we develop a model of indirect reciprocity in which observers revise their private judgements with inertia (Fig.~1). An established reputation changes only after a run of contrary assessments, and the run required to lose standing can differ from the run required to regain it. Negative assessments move an individual in good standing through warning states, and positive assessments move an individual in bad standing through recovery states. Each intermediate state can already carry partial consequences, because it is interpreted as good or bad with an adjustable probability. Observers share no information and update on one observation at a time, so any agreement among them must emerge from how they revise their judgements.

\end{draft}
\begin{draft}{Storyline 3, paragraph 5, 154 words}
% Storyline 3, paragraph 5: findings and implication.
We find that the direction of the revision rule, rather than the amount of evidence it requires, decides whether private observers agree with one another and still distinguish cooperators from defectors. The beneficial direction depends on the norm, mirroring the kind of misjudgement that private assessment produces under each norm. Compact rules, selected solely by these reputational criteria, make cooperation stable and easy to reach under four classical norms, whereas the same amount of inertia placed on the wrong path can leave no cooperation at all. The rule that best protects cooperators from disputed evidence, however, also shelters opportunists who defect only occasionally, a weakness that the aggregation and score-based rules we test share. Allowing warning states to carry partial consequences restores deterrence while cooperation remains high. Sustaining cooperation through private reputations is therefore not only a matter of what individuals know about each other, but also of how they revise what they believe.

\end{draft}

% =============================================================================
\section{Comparison and recommendation}

\begin{center}\small
\begin{tabularx}{\linewidth}{@{}>{\raggedright\arraybackslash}p{2.3cm} L L L@{}}
\toprule
& \textbf{Storyline 1} & \textbf{Storyline 2} & \textbf{Storyline 3} \\
\midrule
Frame & How reputations change & Where agreement comes from & From information to judgement \\
Gap sentence & How direction and the treatment of incomplete evidence affect the dilemma & Which properties of private revision yield accurate agreement & How the revision rule shapes the dilemma, and how it should be designed \\
New concept & Directional inertia, graded expression & Accurate agreement without sharing & Judgement as a trade-off between two errors, managed by direction and partial consequences \\
Closest work & Scores, records, aggregation & Michel-Mata et al.\cite{michelmata2024}, Tham et al.\cite{tham2026} & The whole information-based literature, as a contrast rather than a competitor \\
Fit to Figs.\ 2--5 & Direct for all & Direct for 2--3 & Direct for all, Fig.~5a strongest \\
Breadth & Field & Broad, but borrowed & Broad and distinctive \\
Main risk & Reads as a model extension & Overlap with recent work & Overreach, so the Discussion must carry it \\
Best venue & \textit{PNAS}, \textit{Nat.\ Commun.} & As a supplementary perspective & \textit{Nature}, \textit{Nat.\ Hum.\ Behav.} \\
\bottomrule
\end{tabularx}
\end{center}

\textbf{Recommendation.} For \textit{Nature}, use Storyline~3. It keeps everything you like about Storyline~1, which becomes its core content, and it lifts the frame from a modelling choice to a general problem of social judgement. It also includes Storyline~2 naturally, because agreement without shared information is a consequence of judgement rules. For a field journal, Storyline~1 is the safer choice. In both cases add the Storyline~2 sentence in paragraph~5 and its Discussion paragraph.

\textbf{Matching changes elsewhere.}
\begin{itemize}[leftmargin=1.2em]
\item \textbf{Abstract gap sentence.}
  \begin{itemize}[leftmargin=1.2em]
  \item Storyline~3: ``Existing remedies improve what observers know, but how individuals should revise their judgements as evidence arrives has received little attention.''
  \item Storyline~1: ``Most models let a single assessment change a reputation, or let evidence accumulate symmetrically, and it is unclear how the direction of reputational change and the treatment of incomplete evidence affect cooperation.''
  \end{itemize}
\item \textbf{Title options for Storyline~3.}
  \begin{itemize}[leftmargin=1.2em]
  \item ``How private judgements are revised decides cooperation by reputation''
  \item ``Directional inertia in social judgement sustains cooperation''
  \end{itemize}
\item \textbf{Venue format.} \textit{Nature} articles open with a referenced summary paragraph rather than a conventional abstract. Please check the current guide to authors.
\end{itemize}

\textbf{To check before finalizing.}
\begin{itemize}[leftmargin=1.2em]
\item Yamamoto et al.\ 2025\cite{yamamoto2025}: whether it lets gains and losses differ, or grades treatment. I still could not access it from here.
\item The interpretive sentence in Storyline~3, paragraph~5 (``mirroring the kind of misjudgement \dots''). It rests on the Discussion's argument for the four norms studied. Keep the qualifier if the Discussion changes.
\item Sources for any added examples, such as penalty points or platform strike policies.
\end{itemize}

\bibliographystyle{../../latex_file/naturemag-doi}
\bibliography{../../latex_file/references}
\end{document}
